% ============ §4 形容词 ============
\hdr{§4}{形容词：词尾跟着名词走，拼写跟着词干走}

\callout{\textbf{这一节的坐标}：У7 p26--p34（како́й 四形式、五族词尾）、p36--p37（13 个颜色词）、p39--p42（拼写规则的六组练习）、p44--p46（答 Како́й? 的模型）；У8 p2（指示代词与最高级，见 \S5）。\\
\textbf{这一节要解决的事}：见到一个形容词，能写出它的四个形式；见到 г／к／х／ж／ш／ч／щ 结尾的词干，知道为什么不能写 ы。}

\subhdr{4.1　位置固定：形容词永远在名词前面}

У7 p26 只讲了两件事，第二件是位置：\textbf{形容词总是放在名词前}——\textbf{но́вый телефо́н}（新手机）、\textbf{краси́вая де́вушка}（漂亮姑娘）。倒过来说 \textbf{*телефо́н но́вый} 在陈述句里不合语序。

第一件事是「形容词回答什么问题」，它给出四个疑问词，\textbf{这四个疑问词已经告诉你答案该用哪个词尾}：

\begin{minipage}{\linewidth}
\tabtop
\begin{tabular}{@{}>{\raggedright\arraybackslash}p{0.16\linewidth}>{\raggedright\arraybackslash}p{0.22\linewidth}>{\raggedright\arraybackslash}p{0.26\linewidth}>{\raggedright\arraybackslash}p{0.28\linewidth}@{}}
\thead{疑问词} & \thead{问哪一类名词} & \thead{例问句} & \thead{例答句} \\
\end{tabular}
\tabmid
\begin{tabular}{@{}>{\raggedright\arraybackslash}p{0.16\linewidth}>{\raggedright\arraybackslash}p{0.22\linewidth}>{\raggedright\arraybackslash}p{0.26\linewidth}>{\raggedright\arraybackslash}p{0.28\linewidth}@{}}
Како́й? & 阳性（он） & Како́й э́то дом? & Э́то но́вый дом. \\
Кака́я? & 阴性（она́） & Кака́я э́то у́лица? & Э́то дли́нная у́лица. \\
Како́е? & 中性（оно́） & Како́е э́то фо́то? & Э́то ста́рое фо́то. \\
Каки́е? & 复数（они́） & Каки́е э́то студе́нты? & Э́то у́мные студе́нты. \\
\end{tabular}
\tabbot
\end{minipage}

\textbf{这张表在说什么}：四个疑问词的词尾本身就把性数标出来了——\textbf{-о́й}（阳）、\textbf{-а́я}（阴）、\textbf{-о́е}（中）、\textbf{-и́е}（复）。所以答题时不需要另外判断，\textbf{照抄问句疑问词的词尾}就对了。

\textbf{逐行读法}：第三行 \textbf{Како́е э́то фо́то?} 的答句里 фото́ 是中性不变格词，形容词仍要写中性 -ое（ста́рое）；第四行的 студе́нты 是复数，形容词写 -ые（у́мные）。

\subhdr{4.2　五族词尾总表：一行一个原形词，右边四格是它的四个形式}

У7 p26--p34 与 p41--p42 合成下表——\textbf{全篇最重要的一张}，后面 \S5 的指示代词、\S5.2 的最高级都用同一套词尾。

\textbf{先把「原形」说清}：俄语的形容词在词典、单词表、板书里\textbf{只收一个形式}——\textbf{阳性单数第一格}（回答 \textbf{Како́й?} 的那一格），叫\textbf{原形}；阴性、中性、复数都是它的变形：

\begin{center}
原形 \textbf{но́вый}（阳，он）\ $\to$ \ но́вая（阴，она́）\ $\to$ \ но́вое（中，оно́）\ $\to$ \ но́вые（复，они́）
\end{center}

所以\textbf{原形永远排第一}：看到 но́вая 要知道它属于 но́вый。下表\textbf{一行一个原形词}，右边四格是它的四个形式，\textbf{每格都是完整的词}，可直接查词典、对词表：

\begin{minipage}{\linewidth}
\small
\renewcommand{\arraystretch}{1.10}
\tabtop
\begin{tabular}{@{}>{\raggedright\arraybackslash}p{0.15\linewidth}>{\raggedright\arraybackslash}p{0.15\linewidth}>{\raggedright\arraybackslash}p{0.15\linewidth}>{\raggedright\arraybackslash}p{0.15\linewidth}>{\raggedright\arraybackslash}p{0.17\linewidth}@{}}
\thead{原形（он）} & \thead{阴性（она́）} & \thead{中性（оно́）} & \thead{复数（они́）} & \thead{这一族} \\
\end{tabular}
\tabmid
\begin{tabular}{@{}>{\raggedright\arraybackslash}p{0.15\linewidth}>{\raggedright\arraybackslash}p{0.15\linewidth}>{\raggedright\arraybackslash}p{0.15\linewidth}>{\raggedright\arraybackslash}p{0.15\linewidth}>{\raggedright\arraybackslash}p{0.17\linewidth}@{}}
но́вый & но́вая & но́вое & но́вые & 硬词干 \\
высо́кий & высо́кая & высо́кое & высо́кие & г／к／х \\
хоро́ший & хоро́шая & хоро́шее & хоро́шие & ж／ш／ч／щ \\
си́ний & си́няя & си́нее & си́ние & 软词干 \\
голубо́й & голуба́я & голубо́е & голубы́е & 词尾重音 \\
\end{tabular}
\tabbot
\end{minipage}

\textbf{这张表在说什么}：五行不是五个规则，而是\textbf{同一条规则的五次读数}——\textbf{词干是什么，词尾就跟什么}。硬词干（но́вый）走 -ый 一族；词干末尾是 г／к／х（высо́кий）或 ж／ш／ч／щ（хоро́ший）时，阳性只能写 -ий，其余三格里只有中性换成 -ее；软词干（си́ний）阳性也写 -ий，但阴性／中性／复数换成 -яя／-ее／-ие；词尾重音（голубо́й）另走一套 -о́й／-а́я／-о́е／-ы́е。

\textbf{把词尾单抽出来}（\S5 的指示代词与最高级直接用这一行，看到陌生形容词也用它倒查族别）：

\begin{itemize}
\item \textbf{阳性（он）}：-ый（硬）／-ий（г／к／х、ж／ш／ч／щ、软）／-о́й（重音）
\item \textbf{阴性（она́）}：-ая（硬、г／к／х、ж／ш／ч／щ）／-яя（软）／-а́я（重音）
\item \textbf{中性（оно́）}：-ое（硬、г／к／х）／-ее（ж／ш／ч／щ、软）／-о́е（重音）
\item \textbf{复数（они́）}：-ые（硬）／-ие（г／к／х、ж／ш／ч／щ、软）／-ы́е（重音）
\end{itemize}

\textbf{逐行读法}：五行各是 У7 p28 的原例，顺序固定为\textbf{原形 → 阴 → 中 → 复}（原形放第一位不是排版习惯，是这个词在词典里的位置）：

\begin{enumerate}
\item \textbf{но́вый}：硬词干 но́в-，四个词尾 -ый／-ая／-ое／-ые 原样贴上，没有任何变形。
\item \textbf{высо́кий}：词干末尾是 к，к 后面只能写 и（\S4.3），所以阳性是 -ий，不是 *высо́кый；阴、中、复照常 -ая／-ое／-ие。
\item \textbf{хоро́ший}：词干末尾是 ш，同上写 -ий；中性写成 -ее（хоро́шее）而不是 *хоро́шое，因为 ж／ш／ч／щ 后面不写非重音的 о。
\item \textbf{си́ний}：软词干（си́н- 的 н 是软辅音），阳性也是 -ий，但另外三格换成 -яя／-ее／-ие——这一族不能写 *си́ная。
\item \textbf{голубо́й}：重音落在词尾，整族换成 -о́й／-а́я／-о́е／-ы́е；它的词干末尾是 б（不在 ж／ш／ч／щ 里），所以复数照写 ы（голубы́е），下一段专门说这一格。
\end{enumerate}

\textbf{从形式倒推回原形}（查词典、查本册词表都要走这一步）：阴性 -ая／-яя、中性 -ое／-ее、复数 -ые／-ие 一律把词尾换回 \textbf{-ый／-ий}，就得到原形——кра́сная $\to$ кра́сный，си́нее $\to$ си́ний；-а́я／-о́е／-ы́е 这一族的原形是 \textbf{-о́й}——голуба́я $\to$ голубо́й。

\textbf{表后必须解释的一格}：\textbf{голубы́е 写 ы} 看上去违反了「ж／ш／ч／щ 后写 и」，其实没有——голубо́й 的词干末尾是 \textbf{б}，б 不在这七个字母里，所以复数照写 -ы́е。这条反例在 4.3 会再出现一次。

\aside{「软词干」这一族的关键：词干以软辅音 нь 收尾，阳性词尾不能再写 -ый（硬软冲突），只能写 -ий。си́ний 的 -ий 不是拼写规则要求的，是词干本身软带来的。另外记住\textbf{词干是什么}：把词尾去掉，剩下的就是词干，它是这个形容词自己的骨头——но́вый 的词干是 но́в-，си́ний 的是 си́н-；\textbf{词干不动，动的只有挂在后面的词尾}。}

\needspace{7.4cm}
\subhdr{4.3　拼写规则：г／к／х、ж／ш／ч／щ 后面只写 и，不写 ы}

У7 p41--p42 用两组练习把这条规则讲了六遍。规则本身一句话：

\begin{minipage}{\linewidth}
\tabtop
\begin{tabular}{@{}>{\raggedright\arraybackslash}p{0.24\linewidth}>{\raggedright\arraybackslash}p{0.32\linewidth}>{\raggedright\arraybackslash}p{0.24\linewidth}@{}}
\thead{规则} & \thead{课件例} & \thead{别写成} \\
\end{tabular}
\tabmid
\begin{tabular}{@{}>{\raggedright\arraybackslash}p{0.24\linewidth}>{\raggedright\arraybackslash}p{0.32\linewidth}>{\raggedright\arraybackslash}p{0.24\linewidth}@{}}
词干末尾 г／к／х → 阳性 -ий & ру́сский, ма́ленький, высо́кий, ти́хий & *ру́скый \\
词干末尾 ж／ш／ч／щ → 阳性 -ий & хоро́ший, све́жий, горя́чий, мла́дший & *хоро́шый \\
复数也受同一条约束 & больши́е, плохи́е, ти́хие & *большы́е \\
重音落到词尾 → 换 -о́й 族 & большо́й, чужо́й, дорого́й, плохо́й, молодо́й & — \\
\end{tabular}
\tabbot
\end{minipage}

\textbf{这张表在说什么}：七个字母（г、к、х、ж、ш、ч、щ）在俄语正字法里\textbf{永远不与 ы 相邻}——所以阳性 -ый 碰到它们要让位给 -ий，复数的 -ые 让位给 -ие。这条规则不需要理解，只需要执行。

\textbf{逐行读法}：第三行是这条规则的「覆盖面」——不只是阳性单数，复数也一样；第四行是另一条平行规则：当重音落在词尾，形容词整体换成 -о́й 一族（большо́й、чужо́й），这一族里 \textbf{больши́е 写 и} 又回来了（因为词干末尾是 ш）。

\warnbox{\textbf{边界}：\textbf{голубы́е}（词干末尾 б）与 \textbf{больши́е}（词干末尾 ш）看着像同类，写法不同——判断依据永远是\textbf{词干最后一个字母}，不是整个词看起来像谁。}

\subhdr{4.4　13 个颜色词：一套词尾，三种走法}

У7 p36--p37 把颜色词按词尾分成三型，其实是 4.2 那五族的再一次落地：

\begin{minipage}{\linewidth}
\tabtop
\begin{tabular}{@{}>{\raggedright\arraybackslash}p{0.20\linewidth}>{\raggedright\arraybackslash}p{0.12\linewidth}>{\raggedright\arraybackslash}p{0.20\linewidth}>{\raggedright\arraybackslash}p{0.12\linewidth}>{\raggedright\arraybackslash}p{0.20\linewidth}@{}}
\thead{俄语} & \thead{中文} & \thead{俄语} & \thead{中文} & \thead{词尾型} \\
\end{tabular}
\tabmid
\begin{tabular}{@{}>{\raggedright\arraybackslash}p{0.20\linewidth}>{\raggedright\arraybackslash}p{0.12\linewidth}>{\raggedright\arraybackslash}p{0.20\linewidth}>{\raggedright\arraybackslash}p{0.12\linewidth}>{\raggedright\arraybackslash}p{0.20\linewidth}@{}}
бе́лый & 白 & чёрный & 黑 & -ый 族 \\
кра́сный & 红 & зелёный & 绿 & -ый 族 \\
жёлтый & 黄 & се́рый & 灰 & -ый 族 \\
кори́чневый & 棕 & ора́нжевый & 橙 & -ый 族 \\
ро́зовый & 粉 & фиоле́товый & 紫 & -ый 族 \\
разноцве́тный & 彩色的 & си́ний & 深蓝 & -ий 族（软） \\
голубо́й & 浅蓝 & — & — & -о́й 族（重音） \\
\end{tabular}
\tabbot
\end{minipage}

\textbf{这张表在说什么}：13 个颜色词里，词尾走 -ый 一类的有 10 个，走 -ий 一类的 1 个（си́ний），走 -о́й 一类的 1 个（голубо́й），另有一个 разноцве́тный 也归 -ый 一类。所以「颜色词」不是一类词尾，而是\textbf{三个族都有颜色词}。

\textbf{逐行读法}：这一列「词尾型」直接决定怎么变——\textbf{бе́лый} 走 \textbf{но́вый} 那一行，\textbf{си́ний} 走软词干那一行，\textbf{голубо́й} 走重音在词尾那一行；每一行的四个形式在 4.2 的表里都完整写着。

\textbf{例句}（У7 p37--p39 的模型）：\textbf{кра́сная су́мка}（红包）｜\textbf{си́нее не́бо}（蓝色的天）｜\textbf{чёрный ко́фе}（黑咖啡，ко́фе 是阳性）｜\textbf{голубо́е мо́ре}（浅蓝的海）。

\subhdr{4.5　答 Како́й? 的三步与四个坑}

У7 p44--p46 的模型题全是一个样子：问 \textbf{Како́й э́то дом?}，答 \textbf{Э́то большо́й но́вый дом.}——\textbf{两个形容词都用同一个性数}。做题顺序固定三步：

\begin{enumerate}
\item \textbf{看问句的疑问词}：како́й／кака́я／како́е／каки́е 已经把性数告诉你了。
\item \textbf{选同族词尾}：-ый／-ий／-о́й ↔ -ая／-яя／-а́я ↔ -ое／-ее／-о́е ↔ -ые／-ие／-ы́е。
\item \textbf{形容词放名词前}：不能倒装。
\end{enumerate}

\textbf{四个坑}（У7 与 У8 都踩到过）：

\begin{itemize}
\item \textbf{坑一：不变格外来词}。\textbf{метро́}、\textbf{пальто́}、\textbf{фото́}、\textbf{кафе́} 不变格，性按 2.2 定：\textbf{Како́е э́то метро́?}、\textbf{Э́то моё пальто́}。
\item \textbf{坑二：ко́фе 是阳性}。虽然以 -е 结尾，但它按阳性走：\textbf{Како́й э́то ко́фе?}、\textbf{чёрный ко́фе}。
\item \textbf{坑三：лю́ди 是 челове́к 的复数}。所以问形容词用复数：\textbf{Каки́е лю́ди?}、\textbf{молоды́е лю́ди}。
\item \textbf{坑四：так 与 тако́й}。这一条课件未讲，本讲义只提一句：\textbf{тако́й} 是「这样的」，跟形容词同位置，词尾同五族。
\end{itemize}

\aside{四个坑其实只有一个共同原因：\textbf{形容词的性数必须与名词一致}，而名词的性数不总写在词尾上（外来词、不规则复数）。所以判性顺序永远是：词尾 → 四条 -ь 规律 → 查表（\S2.3）。}

\subhdr{4.6　练习：四个小题}

\begin{enumerate}[label=\textbf{第 \arabic* 题}]
\item （У7 p28）写出 си́ний 的四个形式，并说出它走哪一族。
\item （У7 p40）дома́шний 与 дома́шняя：дома́шний \blank（зада́ние）、дома́шняя \blank（рабо́та）。
\item （У7 p41）改错：*ру́скый → \blank；*хоро́шый → \blank。
\item （У7 p44）答句：— Како́й э́то дом? — Э́то \blank（большой ＋ новый）дом.
\end{enumerate}

\textbf{第 1 题}　关键一步：题目给的就是\textbf{原形} си́ний（阳性单数第一格），它的词干 си́н- 以软辅音 нь 结尾，走 4.2 表里「软词干」那一行。答案：си́ний（原形）／си́няя／си́нее／си́ние。

\textbf{第 2 题}　关键一步：зада́ние 中性（-е）、рабо́та 阴性（-а），形容词跟名词走。答案：дома́шнее ｜ дома́шняя。

\textbf{第 3 题}　关键一步：词干末尾分别是 к 与 ш，都在七个字母里，阳性只能写 -ий。答案：ру́сский ｜ хоро́ший。

\textbf{第 4 题}　关键一步：问句 како́й 已给定阳性，两个形容词都写 -ый／-ий 族。答案：большо́й но́вый（注意 большо́й 是重音在词尾那一族，写 -о́й）。
